%%QCM-PREAMBLE
\documentclass[a4paper]{article}
\usepackage[utf8]{inputenc}
\usepackage[T1]{fontenc}
\usepackage{amsmath, amsfonts, amssymb}
\usepackage{enumitem}
\usepackage[francais,bloc,ensemble,nowatermark]{automultiplechoice}
\usepackage{verbatim,multicol}
\usepackage[most]{tcolorbox}

\newtcolorbox{answerbox}[1][4cm]{
  width=\linewidth, height=#1, colback=white, colframe=black,
  arc=0pt, outer arc=0pt, boxrule=0.4pt,
  left=4pt, right=4pt, top=2pt, bottom=2pt,
}

\def\R{\mathbb{R}}
\def\N{\mathcal{N}}
\def\E{\mathbb{E}}
\def\P{\mathbb{P}}
\def\Z{\mathbb{Z}}

\begin{document}
\AMCrandomseed{1234567}
\def\AMCformQuestion#1{{\sc Question #1 :}}

%%QCM-PREAMBLE-END

%%QCM-EXEMPLAIRE-OPEN ncopies=1
\exemplaire{1}{

%%QCM-HEADER
%%H:establishment{Établissement}
%%H:year{Année 20XX--20XX}
%%H:title{Examen : titre de l'examen}
%%H:duration{Durée : X heures}
%%H:subtitle{Calculatrice autorisée}

\noindent\begin{tabular}{p{0.5\linewidth}p{0.5\linewidth}}
{\sc Établissement} & \hfill Année 20XX--20XX \\
 & \hfill 
\end{tabular}

\begin{center}
\textbf{Examen : titre de l'examen}\\
{\small Durée : X heures}\\
{\small \textit{Calculatrice autorisée}}
\end{center}

%%H:instructions-start
Les réponses sont à donner exclusivement sur la feuille de réponses jointe. Les questions portant le symbole \multiSymbole{} peuvent avoir plusieurs bonnes réponses ; les autres ont une unique bonne réponse.
%%H:instructions-end
%%QCM-HEADER-END

%%QCM-BLOCKS-START
%%QCM-BLOCK bid=t-intro01 kind=text
\section*{Questions}
%%QCM-END bid=t-intro01

%%QCM-BLOCK bid=q-exmult1 kind=question_qcm qtype=mult tag=exemple_multiple
\begin{questionmult}{exemple_multiple}
Quelle(s) affirmation(s) ci-dessous est (sont) correcte(s) ?
\begin{reponses}
\bonne{Une bonne réponse}\bareme{b=1/2,m=0}
\bonne{Une autre bonne réponse}\bareme{b=1/2,m=0}
\mauvaise{Une mauvaise réponse}\bareme{b=0,m=1/2}
\mauvaise{Encore une mauvaise réponse}\bareme{b=0,m=1/2}
\end{reponses}
\end{questionmult}
%%QCM-END bid=q-exmult1

%%QCM-BLOCK bid=q-exsing1 kind=question_qcm qtype=single tag=exemple_simple
\begin{question}{exemple_simple}
Quelle est la bonne réponse ?
\begin{reponses}
\bonne{La bonne réponse}
\mauvaise{Une mauvaise réponse}
\mauvaise{Une autre mauvaise réponse}
\mauvaise{Une dernière mauvaise réponse}
\end{reponses}
\end{question}
%%QCM-END bid=q-exsing1

%%QCM-BLOCK bid=a-justif1 kind=answerbox placement=inline
\noindent \textbf{Question~\ref{q-a-justif1}} — \textbf{Justifie ta réponse}\par\vspace{2pt}
\begin{answerbox}[4cm]\end{answerbox}
\element{bareme}{%
\begingroup\csname AMC@ordretrue\endcsname\csname AMC@inside@digittrue\endcsname\def\AMCchoiceLabel##1{\the\numexpr\value{##1}-1\relax}\def\AMCanswer##1##2{##1\hspace{0.5em}}
\begin{question}{bareme-a-justif1}
\label{q-a-justif1}%
\begin{choicescustom}
  \wrongchoice{0}\scoring{0}
  \correctchoice{1}\scoring{1}
  \correctchoice{2}\scoring{2}
  \correctchoice{3}\scoring{3}
\end{choicescustom}
\end{question}
\endgroup
}
%%QCM-END bid=a-justif1

%%QCM-BLOCKS-END

\newpage
%%QCM-ANSWER-SHEET
%%A:id_grid_digits{4}
%%A:name_field{1}
%%A:columns{2}

\AMCdebutFormulaire

\vspace{0.5cm}\hrule\vspace{0.5cm}
{\Large\bf\begin{center}{Feuille de réponses}\end{center}}

{\noindent \large\bf \underline{Identifiants} :}\\

\vskip5mm {\setlength{\parindent}{0pt}\hspace*{\fill}\AMCcodeGridInt{etu}{4}\hspace*{\fill}
\begin{minipage}[b]{6.5cm}
$\longleftarrow{}$\hspace{0pt plus 1cm} codez votre numéro d'étu\-diant ci-contre, et inscrivez votre nom et prénom ci-dessous $\downarrow\downarrow\downarrow$

\vspace{3ex}

\hfill\champnom{\fbox{
    \begin{minipage}{.9\linewidth}
      Nom et prénom :
      
      \vspace*{.5cm}\dotfill

      \vspace*{.5cm}\dotfill
      \vspace*{1mm}
    \end{minipage}
  }}\hfill\vspace{5ex}\end{minipage}\hspace*{\fill}
}

\vspace{0.7cm}

{\noindent \large\bf \underline{Réponses} :}\\

{  \noindent \bf Les réponses aux questions sont à donner exclusivement ci-dessous ; les réponses données ailleurs ne seront pas prises en compte.
}

\noindent   \underline{Attention} : les cases doivent être noircies et non simplement cochées. En cas d'erreur, utiliser du Tipp-Ex pour effacer complètement la case (NE PAS redessiner la case effacée).

\vspace{0.5cm}

\begin{multicols}{2}\columnseprule=.4pt

\formulaire

\end{multicols}

\ifcsname open@k\endcsname\insertgroup{open}\fi

\ifcsname bareme@k\endcsname
\begin{tcolorbox}[title=R\'eserv\'e correcteur, colback=white, colframe=black!55, fonttitle=\bfseries\small, boxrule=0.5pt, arc=2pt, left=4pt, right=4pt, top=2pt, bottom=2pt, before skip=8pt, after skip=4pt]
\begin{multicols}{2}\columnseprule=.4pt
{\csname AMC@shuffleGfalse\endcsname\insertgroup{bareme}}
\end{multicols}
\end{tcolorbox}
\fi
%%QCM-ANSWER-SHEET-END
}
%%QCM-EXEMPLAIRE-CLOSE
\end{document}
